% =============================================================================
% CHƯƠNG 2 — LƯỢC KHẢO VÀ CƠ SỞ LÝ THUYẾT
% Văn phong khoa học, thuật ngữ và cấu trúc đã được chuẩn hóa.
% =============================================================================

\chapter{LƯỢC KHẢO VÀ CƠ SỞ LÝ THUYẾT}
\label{chap:co-so-ly-thuyet}

\section{Quyền chọn rainbow và bài toán hedging rủi ro}
\label{sec:c2-rainbow}

\subsection{Cấu trúc khoản thanh toán đa tài sản}

Quyền chọn rainbow là nhóm sản phẩm phái sinh có khoản thanh toán phụ thuộc
đồng thời vào từ hai tài sản cơ sở trở lên. Khác với quyền chọn đơn tài sản,
giá trị của quyền chọn rainbow chịu tác động của phân phối biên từng tài sản
và của cấu trúc phụ thuộc chéo, đặc biệt là sự đồng biến trong các giai đoạn
cực đoan \parencite{Johnson1987,Wystup2006}. Khi khoản thanh toán phụ
thuộc vào trung bình theo thời gian, phân phối của toàn bộ quỹ đạo trở thành
đối tượng cần mô hình hóa, thay vì chỉ phân phối giá cuối kỳ.

Với $d$ tài sản, ký hiệu
\begin{equation}
    \bar R_i=\frac{1}{N}\sum_{n=1}^{N}\frac{S_n^i}{S_0^i},
    \qquad i=1,\ldots,d,
    \label{eq:c2-path-average}
\end{equation}
là giá chuẩn hóa trung bình theo thời gian của tài sản $i$ trong $N$ ngày giao
dịch. Hai khoản thanh toán được nghiên cứu là
\begin{align}
    \Phi_{\mathrm{basket}}
    &=\left(\frac{1}{d}\sum_{i=1}^{d}\bar R_i-\kappa\right)^{+},
    \label{eq:c2-basket-payoff}\\
    \Phi_{\mathrm{worst}}
    &=\left(\kappa-\min_{1\le i\le d}\bar R_i\right)^{+},
    \label{eq:c2-worst-payoff}
\end{align}
trong đó $(x)^{+}=\max(x,0)$ và $\kappa$ là mức giá thực hiện chuẩn hóa. Công thức
\eqref{eq:c2-basket-payoff} là quyền chọn mua Asian trên rổ tài sản
(\textit{basket Asian call}); công thức \eqref{eq:c2-worst-payoff} là quyền
chọn bán Asian trên tài sản kém nhất (\textit{Asian worst-of put}).

Hai khoản thanh toán phụ thuộc vào những thành phần khác nhau của phân phối đa biến. Basket Asian
call phụ thuộc vào trung bình chéo nên nhạy với xu hướng chung của rổ. Ngược
lại, worst-of put sử dụng toán tử cực tiểu nên nhạy với tài sản suy giảm mạnh
nhất và với phụ thuộc đuôi trái giữa các tài sản. Phép lấy trung bình theo thời
gian làm giảm tác động của một quan sát riêng lẻ, nhưng đồng thời khiến hợp
đồng phụ thuộc đường đi và làm mất phần lớn công thức đóng vốn có ở quyền chọn
châu Âu \parencite{KemnaVorst1990}.

\subsection{Chi phí giao dịch và thị trường không hoàn chỉnh}

Trong mô hình Black--Scholes lý tưởng, một quyền chọn có thể được sao chép bằng
danh mục tự tài trợ được tái cân bằng liên tục \parencite{BlackScholes1973}.
Với $d$ tài sản, delta trở thành một vector gồm $d$ vị thế. Tuy nhiên, khi mỗi
lần thay đổi vị thế đều phát sinh chi phí tỷ lệ, tổng chi phí có thể tăng nhanh
theo tần suất tái cân bằng; chiến lược sao chép liên tục vì thế không còn khả
thi \parencite{Leland1985}. Đối với khoản thanh toán Asian đa tài sản, khó khăn còn đến
từ việc không có công thức delta đóng và từ sai số mô hình của phân phối đồng
thời.

Do đó, chiến lược giao dịch được tối ưu tại các thời điểm rời rạc để giảm sai
số hedging sau chi phí. Trong thị trường không hoàn chỉnh, không nhất thiết
tồn tại chiến lược làm sai số bằng không trên mọi quỹ đạo. Nghiệm phụ thuộc vào
hàm rủi ro được sử dụng.

\section{Từ hedging rủi ro cổ điển đến \textit{Deep Hedging}}
\label{sec:c2-deep-hedging}

\subsection{Sao chép cổ điển và các giới hạn}

Lý thuyết sao chép cổ điển là chuẩn tham chiếu của bài toán hedging: nếu thị
trường hoàn hảo, động học tài sản đúng mô hình, giao dịch liên tục và không có
chi phí, giá quyền chọn được xác định bởi danh mục sao chép. Các giả định này
không đồng thời thỏa mãn trong thiết kế của luận văn. Hợp đồng phụ thuộc đường
đi, giao dịch theo ngày, có chi phí tỷ lệ, và động học đa tài sản không biết
trước. Vì vậy, sai số hedging là biến ngẫu nhiên không thể loại bỏ hoàn
toàn.

Trong thị trường không hoàn chỉnh, chiến lược có thể được xác định bằng cách
tối thiểu hóa một hàm rủi ro của sai số. Cách tiếp cận mean--variance chọn vị
thế để giảm bình phương sai số \parencite{FollmerSondermann1986,Schweizer2010};
các cách tiếp cận hướng đuôi sử dụng VaR hoặc CVaR để tập trung vào các tổn thất
lớn. Như vậy, mục tiêu chuyển từ ``sao chép chính xác'' sang ``tối ưu chính
sách theo tiêu chuẩn rủi ro đã công bố''.

\subsection{Mạng nơ-ron như một chính sách hedging động}

\textit{Deep Hedging} tham số hóa chính sách giao dịch bằng mạng nơ-ron
\parencite{Buehler2019}. Tại thời điểm $t$, trạng thái $X_t$ được ánh xạ thành
vector vị thế
\begin{equation}
    \delta_t=f_{\theta}(X_t),
    \label{eq:c2-neural-policy}
\end{equation}
trong đó $\theta$ là tập tham số được học từ các quỹ đạo mô phỏng. Hàm mất mát
được tính trên khoản thanh toán, lãi/lỗ của danh mục hedging và chi phí giao dịch của từng quỹ đạo;
gradient được truyền ngược qua toàn bộ chuỗi tái cân bằng. Phương pháp này xử
lý trực tiếp nhiều tài sản, khoản thanh toán phi tuyến, biến trạng thái phụ thuộc đường
đi và chi phí giao dịch mà không cần một công thức định giá đóng.

Chính sách được ước lượng theo phân phối quỹ đạo huấn luyện. Sai lệch về đuôi
phân phối, tương quan theo thời gian hoặc sự chuyển dịch giữa các giai đoạn thị
trường làm thay đổi kết quả hedging trên dữ liệu thực. Do đó, nghiên cứu giữ
cố định kiến trúc mạng và quy trình tối ưu, đồng thời chỉ thay đổi mô hình sinh
dữ liệu.

\section{Thước đo rủi ro và mức vốn ban đầu}
\label{sec:c2-risk-measures}

\subsection{Sai số bình phương trung bình}

Ký hiệu $R=\Phi-V_T$ là phần nghĩa vụ chưa được bù đắp tại đáo hạn. Hàm mất
mát bình phương là
\begin{equation}
    \rho_{\mathrm{MSE}}(R)=\mathbb{E}[R^2].
    \label{eq:c2-mse}
\end{equation}
Mọi quan sát trong mẫu đều đóng góp vào gradient của MSE. Hàm mất mát này xử
lý đối xứng sai số dương và âm. MSE được dùng để ước lượng chiến lược theo tiêu
chuẩn bình phương sai số và xác định $V_0^{*}$.

\subsection{Value-at-Risk và Conditional Value-at-Risk}

$\mathrm{VaR}_{\alpha}(R)$ là phân vị mức $\alpha$ của phân phối sai số.
$\mathrm{CVaR}_{\alpha}(R)$ là giá trị trung bình của sai số trong phần đuôi
bất lợi vượt ngưỡng VaR. Theo biểu diễn Rockafellar--Uryasev
\parencite{Rockafellar2000},
\begin{equation}
    \mathrm{CVaR}_{\alpha}(R)
    =\min_{\nu\in\mathbb{R}}
    \left\{\nu+\frac{1}{1-\alpha}
    \mathbb{E}\left[(R-\nu)^{+}\right]\right\}.
    \label{eq:c2-ru}
\end{equation}
Với $\alpha=0{,}95$, gradient tập trung vào khoảng $5\%$ quan sát có sai số lớn
nhất. Biểu diễn \eqref{eq:c2-ru} biến CVaR thành một mục tiêu có thể tối ưu bằng
mini-batch thông qua tham số phụ $\nu$.

\subsection{Mức vốn ban đầu $V_0^{*}$}

Trong giai đoạn MSE, số tiền ban đầu $V_0$ được học cùng chính sách. Nếu
\begin{equation}
    R=\Phi-V_0-\mathrm{PnL}(\delta)+\mathrm{cost}(\delta),
\end{equation}
thì điều kiện bậc nhất theo $V_0$ cho
\begin{equation}
    \frac{\partial}{\partial V_0}\mathbb{E}[R^2]
    =-2\mathbb{E}[R]=0.
    \label{eq:c2-v0-foc}
\end{equation}
Do đó, $V_0^{*}$ làm sai số kỳ vọng bằng không dưới phân phối huấn luyện. Đại
lượng này được gọi là giá bàng quan bậc hai (\textit{quadratic indifference
price}) trong tiêu chuẩn trung bình--phương sai
\parencite{FollmerSondermann1986,Schweizer2010}. Trong nghiên cứu này,
$V_0^{*}$ được sử dụng như mức vốn ban đầu cần thiết để thực hiện chiến lược.
Đại lượng này phụ thuộc vào mô hình sinh dữ liệu, chiến lược và chi phí giao
dịch; nó không phải giá thị trường hoặc giá phi arbitrage duy nhất.

CVaR có tính tịnh tiến:
$\mathrm{CVaR}_{\alpha}(X-a)=\mathrm{CVaR}_{\alpha}(X)-a$. Nếu tiếp tục cho
$V_0$ học tự do trong giai đoạn CVaR, đạo hàm theo $V_0$ là hằng số và bài toán
không có nghiệm dừng hữu hạn. Vì vậy, quy trình thực nghiệm xác định $V_0^{*}$
ở giai đoạn MSE rồi cố định tham số này khi tinh chỉnh CVaR.

\section{Các mô hình sinh dữ liệu tài chính đa biến}
\label{sec:c2-generators}

\subsection{Chuyển động Brown hình học đa biến}

GBM đa biến giả định
\begin{equation}
    \frac{\mathrm{d}S_t^i}{S_t^i}
    =\mu_i\,\mathrm{d}t+\sigma_i\,\mathrm{d}W_t^i,
    \qquad
    \mathrm{Corr}(\mathrm{d}W_t^i,\mathrm{d}W_t^j)=\rho_{ij}.
    \label{eq:c2-gbm}
\end{equation}
GBM được hiệu chỉnh trực tiếp từ lợi suất trung bình, độ biến động và ma trận
tương quan mẫu. Giả định lợi suất log Gaussian cùng biến động và tương quan
không đổi không mô tả đuôi dày, bất đối xứng và phụ thuộc đuôi của dữ liệu cổ
phiếu.

\subsection{Mô hình Heston đa biến}

Heston đưa thêm phương sai ngẫu nhiên vào từng tài sản
\parencite{Heston1993}:
\begin{align}
    \mathrm{d}\log S_t^i
    &=\left(\mu_i-\tfrac12 v_t^i\right)\mathrm{d}t
      +\sqrt{v_t^i}\,\mathrm{d}W_t^{S,i},\\
    \mathrm{d}v_t^i
    &=\kappa_v(\vartheta_i-v_t^i)\mathrm{d}t
      +\xi\sqrt{v_t^i}\,\mathrm{d}W_t^{v,i}.
    \label{eq:c2-heston}
\end{align}
Tương quan âm giữa cú sốc giá và phương sai tạo hiệu ứng đòn bẩy. Phương sai
ngẫu nhiên tạo cụm biến động và đuôi dày hơn GBM. Trong trường hợp đa tài sản,
đặc tả giản lược sử dụng chung $\kappa_v$, $\xi$ và tương quan đòn bẩy giữa các
tài sản để giới hạn số tham số \parencite{Dimitroff2011}. Đặc tả này không mô
hình hóa riêng động học phương sai của từng tài sản.

\subsection{Các mô hình sinh dựa trên dữ liệu}

Các mô hình sinh dựa trên dữ liệu ước lượng phân phối và cấu trúc phụ thuộc từ
mẫu lịch sử. Kết quả ước lượng phụ thuộc vào kích thước mẫu, phương pháp làm
trơn và miền hỗ trợ của dữ liệu. SBTS được xây dựng từ bài toán entropy trên
không gian quỹ đạo và được triển khai trên dữ liệu hữu hạn bằng ước lượng
kernel \parencite{Hamdouche2023}.

\section{Nền tảng lý thuyết của \textit{Schrödinger Bridge}}
\label{sec:c2-sb-theory}

\subsection{Bài toán entropy trên không gian quỹ đạo}

Cho $\Omega=C([0,T],\mathbb{R}^{d})$ là không gian các quỹ đạo liên tục và
$\mathbb{P}_{W}$ là độ đo Wiener tham chiếu. Với một độ đo quỹ đạo
$\mathbb{Q}\ll\mathbb{P}_{W}$, độ phân kỳ Kullback--Leibler là
\begin{equation}
    \mathrm{KL}(\mathbb{Q}\|\mathbb{P}_{W})
    =\mathbb{E}^{\mathbb{Q}}
    \left[\log\frac{\mathrm{d}\mathbb{Q}}
    {\mathrm{d}\mathbb{P}_{W}}\right].
    \label{eq:c2-path-kl}
\end{equation}
Bài toán \textit{Schrödinger Bridge} hai biên chọn độ đo quỹ đạo gần độ đo tham chiếu nhất
theo KL nhưng đồng thời khớp phân phối tại đầu và cuối kỳ
\parencite{Leonard2014}:
\begin{equation}
    \mathbb{Q}^{*}
    =\arg\min_{\mathbb{Q}}
    \mathrm{KL}(\mathbb{Q}\|\mathbb{P}_{W})
    \quad\text{với}\quad
    \mathbb{Q}_{0}=\mu_0,\quad \mathbb{Q}_{T}=\mu_T.
    \label{eq:c2-classical-sb}
\end{equation}
Bài toán này liên hệ với vận tải tối ưu có điều chuẩn entropy và điều khiển
ngẫu nhiên \parencite{ChenGeorgiouPavon2021}.

Đối với chuỗi thời gian, ràng buộc chỉ ở hai đầu chưa đủ. Mở rộng nhiều biên
đòi hỏi
\begin{equation}
    \mathbb{Q}_{t_j}=\mu_{t_j},
    \qquad j=0,1,\ldots,N,
    \label{eq:c2-multimarginal-sb}
\end{equation}
để độ đo quỹ đạo nhất quán với chuỗi phân phối quan sát tại nhiều thời điểm.
Hamdouche và cộng sự chuyển ý tưởng này thành mô hình sinh chuỗi thời gian,
trong đó chuyển tiếp kế tiếp được điều kiện hóa bởi phần lịch sử gần nhất
\parencite{Hamdouche2023}.

\subsection{Từ HJB đến phương trình nhiệt bằng Cole--Hopf}

Trên một khoảng $[t_j,t_{j+1})$, biểu diễn điều khiển ngẫu nhiên của bài toán
KL dẫn đến một hàm giá trị $V(t,x)$ thỏa phương trình Hamilton--Jacobi--Bellman
\begin{equation}
    \partial_t V+\frac12\Delta V-\frac12\|\nabla V\|^2=0.
    \label{eq:c2-hjb}
\end{equation}
Phương trình \eqref{eq:c2-hjb} phi tuyến do hạng tử bình phương gradient. Đặt
\begin{equation}
    \psi(t,x)=\exp[-V(t,x)],
    \label{eq:c2-cole-hopf}
\end{equation}
ta có
\begin{align*}
    \partial_t V&=-\frac{\partial_t\psi}{\psi},
    &\nabla V&=-\frac{\nabla\psi}{\psi},
    &\Delta V&=-\frac{\Delta\psi}{\psi}
      +\frac{\|\nabla\psi\|^2}{\psi^2}.
\end{align*}
Thế vào \eqref{eq:c2-hjb}, hai hạng tử bậc hai triệt tiêu và thu được phương
trình nhiệt lùi tuyến tính
\begin{equation}
    \partial_t\psi+\frac12\Delta\psi=0.
    \label{eq:c2-backward-heat}
\end{equation}
Phép biến đổi Cole--Hopf chuyển phương trình HJB phi tuyến thành phương trình
nhiệt tuyến tính. Nghiệm của phương trình tuyến tính được biểu diễn bằng nhân
nhiệt Brown. Chuyển tiếp giữa hai thời điểm quan sát sau đó được ước lượng từ
các điểm cuối lịch sử bằng trọng số kernel.

\subsection{Từ nghiệm tuyến tính đến quy tắc chuyển tiếp thực nghiệm}

Gọi $p_W(t,x;t_{j+1},y)$ là mật độ chuyển tiếp Brown. Nghiệm của
\eqref{eq:c2-backward-heat} với điều kiện cuối $\rho_{j+1}$ có dạng
\begin{equation}
    \psi(t,x)
    =\int p_W(t,x;t_{j+1},y)\rho_{j+1}(y\mid\mathcal{H}_j)\,\mathrm{d}y,
    \label{eq:c2-heat-solution}
\end{equation}
trong đó $\mathcal{H}_j$ là thông tin quá khứ tại thời điểm $t_j$. Trên mẫu hữu
hạn, phân phối có điều kiện trong \eqref{eq:c2-heat-solution} được thay bằng
phân phối thực nghiệm của các quỹ đạo lịch sử gần trạng thái hiện tại. Vì vậy,
việc triển khai không cần giải PDE trên một lưới không gian nhiều chiều; bài
toán trở thành tính trọng số kernel và trung bình các chuyển tiếp quan sát.

\section{Mô hình SBTS thực nghiệm}
\label{sec:c2-sbts}

\subsection{Xấp xỉ Nadaraya--Watson}

Giả sử có $M$ quỹ đạo tham chiếu
$\{\mathbf{X}^{(m)}\}_{m=1}^{M}$. Kỳ vọng có điều kiện của một hàm $g$ được
xấp xỉ bằng bộ ước lượng Nadaraya--Watson
\parencite{Nadaraya1964,Watson1964}:
\begin{equation}
    \widehat{\mathbb{E}}[g(X_{t_{j+1}})\mid\mathcal{H}_j]
    =\frac{\sum_{m=1}^{M}W_{K}^{(m)}
    g(X_{t_{j+1}}^{(m)})}
    {\sum_{m=1}^{M}W_{K}^{(m)}}.
    \label{eq:c2-nw}
\end{equation}
Trọng số được xây dựng từ tích các kernel trên $K$ quan sát gần nhất:
\begin{equation}
    W_{K}^{(m)}
    =\prod_{\ell=\max(0,j-K+1)}^{j}
    \mathcal{K}_{h}
    \left(\mathbf{x}_{\ell}-\mathbf{X}_{t_{\ell}}^{(m)}\right),
    \qquad
    \mathcal{K}_{h}(u)
    =(h^2-\|u\|^2)^2\mathbf{1}_{\{\|u\|<h\}}.
    \label{eq:c2-kernel}
\end{equation}
Băng thông (\textit{bandwidth}) $h$ kiểm soát bán kính lân cận, còn $K$ kiểm soát chiều dài bộ nhớ.
$h$ quá nhỏ hoặc $K$ quá lớn làm số quỹ đạo có trọng số khác không giảm mạnh;
$h$ quá lớn hoặc $K$ quá nhỏ làm ước lượng quá trơn và mất thông tin điều kiện.

Trong thuật toán SBTS, trọng số lịch sử ở \eqref{eq:c2-kernel} được kết hợp
với nhân Brown xuất phát từ nghiệm Cole--Hopf. Kết quả là một quy tắc chuyển
tiếp lấy trung bình các dịch chuyển kế tiếp của quỹ đạo tham chiếu, với trọng
số cao hơn cho những quỹ đạo vừa giống lịch sử gần đây vừa phù hợp với trạng
thái nội suy hiện tại. Sau đó, nhiễu Gaussian ở các bước Euler nhỏ được cộng
vào để tạo ra quỹ đạo mới. Chương 3 trình bày cấu hình số cụ thể dùng trong
thực nghiệm.

\subsection{Mở rộng đa biến và lựa chọn siêu tham số}

Mở rộng đa biến được thực hiện bằng cách thay khoảng cách tuyệt đối một chiều
bằng chuẩn Euclid của vector lợi suất $d$ chiều. Cùng một kernel vì thế đồng
thời đánh giá độ gần của toàn bộ rổ tài sản. Trong nghiên cứu, $(h,K)$ được
chọn theo MSE dự báo trên phần dữ liệu giữ lại:
\begin{equation}
    (h^{*},K^{*})
    =\arg\min_{(h,K)\in\mathcal{G}_{h}\times\mathcal{G}_{K}}
      \mathrm{MSE}_{\mathrm{holdout}}(h,K).
    \label{eq:c2-hk-selection}
\end{equation}
Quy trình lựa chọn bằng dữ liệu giúp giảm tính tùy ý nhưng không loại bỏ lời
nguyền số chiều: khi $d$ hoặc $K$ tăng, lân cận kernel trở nên thưa và nhu cầu
số quỹ đạo tham chiếu tăng nhanh.

\subsection{Khả năng nội suy và giới hạn dưới dịch chuyển phân phối}

Kernel có hỗ trợ hữu hạn nên SBTS thực hiện nội suy từ các cửa sổ lịch sử. Khi
trạng thái đánh giá cách xa tập tham chiếu, số quỹ đạo có trọng số dương giảm
và phương sai ước lượng tăng. Chương 4 kiểm tra thứ hạng của ba mô hình sinh
dữ liệu tại giai đoạn đại diện, phục hồi và căng thẳng. Nghiên cứu không thực
hiện kiểm định nhân quả riêng cho cơ chế suy giảm ngoài miền dữ liệu.

\section{Lược khảo nghiên cứu thực nghiệm liên quan}
\label{sec:c2-empirical-review}

\textit{Deep Hedging} tối ưu trực tiếp chính sách giao dịch trong điều kiện thị
trường không hoàn chỉnh, khoản thanh toán phi tuyến và chi phí giao dịch
\parencite{Buehler2019}. Các nghiên cứu tiếp theo mở rộng hàm rủi ro, mô hình
biến động và kiến trúc mạng, nhưng phần lớn tập trung vào quyền chọn đơn tài
sản. Bằng chứng về \textit{Deep Hedging} cho quyền chọn Asian rainbow với mô
hình sinh dữ liệu đa tài sản còn hạn chế.

SBTS tạo chuỗi mới từ phân phối và chuyển tiếp thực nghiệm
\parencite{Hamdouche2023}. Độ phù hợp thống kê của dữ liệu tổng hợp và kết quả
hedging rủi ro là hai đối tượng đánh giá khác nhau. Nghiên cứu này kiểm tra
trực tiếp sai số của chiến lược trên dữ liệu thị trường, bên cạnh các chỉ tiêu
phân phối.

Về suy luận, kiểm định từng cặp trả lời liệu hai mô hình sinh dữ liệu có khác nhau trong một
tổ hợp thực nghiệm, trong khi Model Confidence Set (MCS) lựa chọn một tập mô hình
không thể phân biệt với mô hình tốt nhất khi nhiều ứng viên được so sánh đồng
thời \parencite{HansenLundeNason2011}. Kết hợp hai cách nhìn giúp tránh kết
luận dựa duy nhất vào một thứ hạng trung bình có thể không ổn định qua seed.

\section{Quy trình phân tích của nghiên cứu}
\label{sec:c2-framework}

Từ cơ sở trên, nghiên cứu tổ chức quan hệ phân tích theo chuỗi sau:
\begin{enumerate}[label=(\arabic*)]
    \item dữ liệu lịch sử xác định tham số hoặc phân phối thực nghiệm của từng
    mô hình sinh dữ liệu;
    \item mỗi mô hình sinh dữ liệu tạo tập quỹ đạo huấn luyện có cùng kích thước và thời
    hạn;
    \item cùng một kiến trúc \textit{Deep Hedging} được huấn luyện trên từng tập quỹ
    đạo bằng quy trình MSE--CVaR;
    \item các checkpoint được đánh giá trên cùng các cửa sổ lịch sử tại ba
    giai đoạn thị trường;
    \item hiệu quả được so sánh bằng độ lệch chuẩn sai số, CVaR, $V_0^{*}$,
    kiểm định ghép cặp có điều chỉnh nhiều phép thử và MCS.
\end{enumerate}
Ba lớp kết quả gồm độ phù hợp phân phối của mô hình sinh dữ liệu, hiệu quả
hedging và mức ổn định giữa các giai đoạn thị trường. Kết luận về hiệu quả
được xác định từ sai số hedging trên dữ liệu đánh giá. Các mô-men của dữ liệu
tổng hợp chỉ được dùng để đánh giá độ phù hợp phân phối.

\section{Tóm tắt chương}

Chương 2 trình bày hai quyền chọn Asian rainbow, \textit{Deep Hedging}, MSE,
CVaR và $V_0^{*}$. Ba mô hình sinh dữ liệu gồm GBM, Heston và SBTS. Đối với
SBTS, bài toán KL nhiều biên được tuyến tính hóa bằng phép biến đổi Cole--Hopf
và được ước lượng bằng phương pháp Nadaraya--Watson đa biến. Khả năng nội suy
của kernel là cơ sở để đánh giá kết quả tại ba giai đoạn thị trường trong các
chương tiếp theo.
